\chapter{Twisted projections and the final contradiction}
\label{chap:twisted}

\section{The slope cannot remain small}

\begin{lemma}[A lower bound for the slope; WZ2, Lemma 6.6]
\label{lem:large-slope}
\lean{Kakeya.Assouad.LargeSlopeTechnicalStatement}
\uses{prop:c2-grains,def:cwa}
For each $\eps>0$, choose $\eta>0$ and the fine scale sufficiently
small. Let a grain configuration satisfy Proposition~\ref{prop:c2-grains}
with loss $\eta$. Suppose $J\subset[-1,1]$ satisfies
$\sqrt\delta\leq |J|\leq\delta^\eps$,
\[
 \sum_{T\in\T}|Y(T)\cap(\R^2\times J)|\geq\delta^\eta|J|,
\]
and $E_Y\cap(\R^2\times J)$ is covered by at most
$\delta^{-\eta}|J|^{-2+\sigma}$ cubes of side $|J|$.
Then $|f'(z)|\geq |J|$ for $z\in J$.
\end{lemma}
\begin{proof}
\uses{prop:sticky-structure}
If $|f'|$ is smaller than $|J|$ somewhere, the $C^2$ bound keeps
$f$ nearly constant across $J$. Global grains and the spatial covering
hypothesis then concentrate many shaded tube pieces in a thin prism.
The balanced-cover and multiplicity bounds turn shaded mass in that
prism into a count of whole tubes in a slightly enlarged convex body.
Its volume is too small for that count to satisfy CWA. The
formalized argument retains the selected shading, its per-tube mass,
and the actual enclosing prism throughout this step.
\end{proof}

\begin{proposition}[Nonsingular smooth grains; WZ2, Proposition 6.5]
\label{prop:large-slope}
\lean{Kakeya.Assouad.PureWZ2LargeSlopeFromCriticalStatement,Kakeya.Assouad.pure_wz2_node06_large_slope}
\leanok
\uses{def:extremal,def:ad-set,def:plane-map}
For every $\eps,\delta_0>0$, there is an $\eps$-extremal
configuration at some $0<\delta\leq\delta_0$ with the local grain
property and a global grain slope $f\in C^2([-1,1])$ satisfying
\[
       1\leq|f'(z)|\leq2,\qquad |f''(z)|\leq1/100
       \quad(-1\leq z\leq1).
\]
Every horizontal slice has the global AD bound of
Proposition~\ref{prop:c2-grains}, with this same $f$ and the same
shading.
\end{proposition}
\begin{proof}
\uses{prop:c2-grains,lem:large-slope,lem:mild-rescaling}
\leanok
Choose a populated interval of heights whose spatial covering number
has the required extremal bound. The preceding lemma gives a lower
bound for the derivative there. Localize further so that the derivative
has comparable magnitude, then rescale the height and transverse
coordinates. A sufficiently short interval makes the second derivative
small after normalization. Transport the family, shading, and both
grain structures together through this rescaling.
\end{proof}

\begin{definition}[Twisted projection; WZ1, Definition 6.1]
\label{def:twisted-projection}
\lean{Kakeya.Assouad.twistedProjection,Kakeya.Assouad.twistedCurve}
\leanok
For $f:[-1,1]\to\R$, define
\[
 \pi_f(x,y,z)=(x+f(z)y,z).
\]
In a vertical coordinate chart, a line has the form
$\ell_{a,b,c,d}(t)=(a+ct,b+dt,t)$, and its image is the graph of
\[
                   t\longmapsto a+ct+f(t)(b+dt).
\]
\end{definition}

\begin{lemma}[A small twisted projection; WZ2, Lemma 6.7, corrected exponent]
\label{lem:small-twisted-projection}
\lean{Kakeya.Assouad.PureWZ2SmallTwistedProjectionData,Kakeya.Assouad.pure_wz2_node07_small_twisted_projection}
\leanok
\uses{def:twisted-projection,def:cwa,def:shading}
Assume a positive critical exponent $\sigma$ exists. For every
$a>0$ and every scale ceiling, there is a nonempty configuration
below that ceiling, with CWA error $\delta^{-a}$, shading density
$\delta^a$, and a nonsingular slope as above, such that
\[
                        |\pi_f(E_Y)|\leq\delta^{\sigma-a}.
\]
The final preparation can also supply, on the same selected
configuration, the parameter-space hypotheses of
Theorem~\ref{thm:twisted-parameter}.
\end{lemma}
\begin{proof}
\uses{prop:large-slope}
\leanok
At each height, the global AD bound covers the horizontal projection
by at most $O(\delta^{-(1-\sigma)-a})$ intervals of length $\delta$.
The cubical, slabwise form of this statement bounds planar area by
$O(\delta^{\sigma-a})$. Allocate a smaller source loss to absorb
the fixed constants. The further analytic preparation retains this
upper bound while producing separated representatives and the
required indexed parameter counts.
\end{proof}

\begin{remark}[The sign in WZ2 Lemma 6.7]
\label{rem:projection-sign}
WZ2 prints $\delta^{\sigma+\eps}$ in this lemma. The grain estimate
with covering loss $\delta^{-\eps}$ instead gives
$\delta^{\sigma-\eps}$, and this is the exponent in
\texttt{PureWZ2SmallTwistedProjectionData}. The corresponding WZ1
Corollary 6.5 is formulated as an $L^{3/2}$ lower bound; the volume
and covering losses in its derivation also have the latter sign.
The plus sign would be a stronger assertion when $0<\delta<1$.
The weaker, minus-sign estimate is sufficient: choose the lower-bound
exponent strictly below $\sigma-a$ in the final contradiction.
\end{remark}

\section{The projection lower bound}

\begin{proposition}[Convex Wolff version; WZ2, Proposition 6.8]
\label{prop:twisted-cwa}
\uses{def:cwa,def:shading,def:twisted-projection}
For every $\eps>0$, the paper asserts that there are
$\eta,\delta_0>0$ such that a nonempty $\delta$-tube configuration
with CWA error $\delta^{-\eta}$ and a $\delta^\eta$-dense shading
satisfies
\[
                         |\pi_f(E_Y)|\geq\delta^\eps
\]
for $0<\delta\leq\delta_0$ and every slope with
$1\leq|f'|\leq2$, $|f''|\leq1/100$ on the height window.
\end{proposition}
\begin{proof}
\uses{prop:twisted-tube-wolff}
A radius-$r$ tube has volume $O(r^2)$. Applying CWA to that convex
body gives the dimension-$2$ Tube Wolff count, with a fixed change
of constant. Apply the paper's Tube Wolff projection estimate with
a smaller source loss to absorb this constant.
\end{proof}

\begin{proposition}[Tube Wolff version; WZ2, Proposition 7.1]
\label{prop:twisted-tube-wolff}
\uses{def:tube-frostman,def:shading,def:twisted-projection}
The paper gives the same conclusion assuming the dimension-$2$
Tube Wolff condition in place of CWA. In a bounded line chart the
intended hypothesis controls the number of line parameters in every
ball of radius $r\geq\delta$ by
$\delta^{-\eta}r^2\#\T$.
\end{proposition}

\begin{remark}[The lower bound used by the formal proof]
\label{rem:parameter-frostman}
WZ2 Remark 7.4 identifies the Tube Wolff condition with Frostman
nonconcentration in the space of line parameters. For the repository's
ordinary unit-segment carriers these are distinct predicates:
containment of a complete segment in a coarse tube does not, by itself,
count every segment with nearby supporting-line parameters. Segments
on the same line can have different longitudinal positions.
The proof therefore supplies the indexed four-parameter Frostman
bound explicitly. It also supplies a bounded chart and a cardinality
upper bound. The next theorem is the actual proved specialization;
the two preceding paper formulations are not marked as independently
proved in their greater generality.
\end{remark}

\begin{theorem}[The parameter-Frostman specialization of WZ2, Proposition 7.1]
\label{thm:twisted-parameter}
\lean{Kakeya.Assouad.TwistedProjectionParameterFrostmanEstimateStatement,Kakeya.Assouad.pure_wz2_twisted_projection_parameter_frostman_lower_bound}
\leanok
\uses{def:twisted-projection,def:tube-frostman,def:shading}
For every $\eps>0$, there are $0<\eta\leq\eps$ and
$0<\delta_0<1$ with the following property. Let $\T$ be a nonempty
family at $0<\delta\leq\delta_0$, represented in a vertical chart
with tube bases of norm at most $4$ and vertical direction component
of absolute value at least $1/2$. Assume analytic essential
distinctness, the dimension-$2$ Tube Wolff bound with error
$\delta^{-\eta}$, and the indexed parameter bound
\[
 \#\{T:\|(a_T,b_T,c_T,d_T)-(a_{T_0},b_{T_0},c_{T_0},d_{T_0})\|_\infty
                  \leq r\}
       \leq\delta^{-\eta}r^2\#\T
       \quad(\delta\leq r\leq1).
\]
Assume also $\#\T\leq\delta^{-2-\eta/10}$.
If $Y$ is $\delta^\eta$-dense and supported in $-1\leq z\leq1$,
and $f(0)=0$, $1\leq|f'|\leq2$, $|f''|\leq1/100$, then
\[
                         |\pi_f(E_Y)|\geq\delta^\eps.
\]
\end{theorem}
\begin{proof}
\uses{lem:spacing,lem:block-replication,lem:projection-one-step,
  thm:cinematic-l32,thm:os-discretized,lem:uniform-parameter-coarsening,
  lem:projection-iteration}
\leanok
Separate the parameters into blocks and regularize their branching
over a finite sequence of scales. The spacing argument finds a
block with sufficient local nonconcentration. Replicate this block
and use the cinematic estimate to obtain the one-step comparison of
fine projection area with coarse projection area. If that comparison
does not yet give the desired lower bound, pass to a coarser scale
and repeat. The scale increases by a controlled power at each step,
so the procedure stops after a uniformly bounded number of steps.
The successive area comparisons telescope. The scale schedule and
the source loss are chosen so that the total loss is less than
$\eps$.

The iteration keeps the induced shading on the actual selected
family. Clipping a thickened shading back to the height window costs
an additive boundary term; this term is estimated before converting
the mass bound into a density bound.
Lemmas~\ref{lem:uniform-parameter-coarsening} and
\ref{lem:projection-iteration} isolate the counting and arithmetic
steps in this passage; they do not assert that an arbitrary coarse
shading has the required density or projection compatibility.
\end{proof}

\section{Discrete sets and scale selection}

\begin{definition}[Katz--Tao and Frostman sets; WZ2, Definitions 7.2--7.3]
\label{def:kt-frostman}
\lean{Kakeya.Assouad.DiscreteSet.IsKatzTao,Kakeya.Assouad.DiscreteSet.IsFrostman}
\leanok
A nonempty bounded set $A\subset\R^n$ is a
$(\delta,s,C)$-Katz--Tao set if
\[
       \cE_\delta(A\cap B(x,r))\leq C(r/\delta)^s
       \quad(\delta\leq r\leq1).
\]
It is a $(\delta,s,C)$-Frostman set if the right-hand side is replaced
by $Cr^s\cE_\delta(A)$. For separated finite sets the formalization
uses point counts; they are comparable to covering numbers with a
dimensional change of constant.
\end{definition}

\begin{remark}[A reversed quotient in Definition 7.2]
\label{rem:kt-quotient}
The April 2025 WZ2 manuscript prints $C(\delta/r)^s$. The inequality
used in its arguments, and in the Lean predicate, has $C(r/\delta)^s$.
For example, a separated set with $\asymp\delta^{-s}$ points in a
unit ball cannot satisfy the printed formula at $r=1$ with bounded
$C$. All Katz--Tao statements here use the latter quotient.
\end{remark}

\begin{lemma}[Frostman to Katz--Tao; WZ2, Lemma 7.6]
\label{lem:frostman-to-kt}
\lean{Kakeya.Assouad.FrostmanToKatzTaoStatement,Kakeya.Assouad.frostman_to_katz_tao}
\lean{Kakeya.Assouad.FrostmanToKatzTaoWeightedCardStatement,Kakeya.Assouad.frostman_to_katz_tao_weighted_card}
\uses{def:kt-frostman}
A $(\delta,s,C)$-Frostman set in a bounded region has a
$(\delta,s,100)$-Katz--Tao subset $A'$ with
\[
                \cE_\delta(A')\gtrsim_{n,s}
                   C^{-1}\delta^{-s}(1+\log(1/\delta))^{-1}.
\]
The linked finite-set form assumes $\delta$-separation and
$\#A\geq\delta^{-s}$ and writes the implicit dimensional constant
explicitly as an existential constant.
The sticky proof uses the linked weighted-cardinality version instead:
it removes the separation hypothesis, assumes
$\delta^{-s}\leq C\#A$, and retains
$\delta^{-s}\leq L(1+\log(\delta^{-1}))C\#A'$.
The dependency audit distinguishes these two Lean versions.
\end{lemma}
\begin{proof}
\uses{lem:content-to-kt}
Use the Frostman bound to control the cost of coverings at each
dyadic scale. A logarithmic pigeonhole produces a portion with
enough finite-scale content. The content extraction lemma then
gives a separated Katz--Tao set of the claimed size.
\end{proof}

\begin{lemma}[Content extraction; WZ2, Lemma 7.7]
\label{lem:content-to-kt}
\uses{def:kt-frostman}
If $\cH^s_\infty(A)=\kappa>0$, there is a
$(\delta,s,100)$-Katz--Tao subset $A'\subset A$ with
$\cE_\delta(A')\gtrsim\kappa\delta^{-s}$.
\end{lemma}
\begin{proof}
Discretize a Frostman measure at scale $\delta$, retaining at most a
fixed amount of mass in each occupied cell. Select representatives
by the usual dyadic pruning procedure. The mass distribution bound
controls the number selected in every larger ball, and the total
mass supplies the lower bound.
\end{proof}

\begin{definition}[Uniform branching; WZ2, Definition 7.8]
\label{def:uniform-set}
Fix a dyadic sequence $1=\Delta_0>\cdots>\Delta_m=\delta$.
A set is uniform on these scales if the number of occupied
$\Delta_{j+1}$-children is the same for every occupied
$\Delta_j$-cube. Consequently the number of finest-scale cells in
each occupied cube at a fixed level is the same. WZ2's
$\eta$-uniformity allows comparability by a factor of $100$ and uses
the scales $2^{-j/\eta}$.
\end{definition}

\begin{lemma}[Uniform refinement; WZ2, Lemma 7.9, in block notation]
\label{lem:uniform-set}
\uses{def:uniform-set}
Let $\delta=2^{-mT}$ and $\Delta_j=2^{-jT}$, $0\leq j\leq m$.
A set $A\subset[0,1)^n$ has a refinement $A'$ uniform on these
scales and satisfying
\[
                   \cE_\delta(A')\geq (C_nT)^{-m}\cE_\delta(A).
\]
In particular, choosing $T$ large enough in terms of $\eps$ gives
$\cE_\delta(A')\geq\delta^\eps\cE_\delta(A)$.
\end{lemma}
\begin{proof}
At one level there are $O_n(T)$ dyadic possibilities for the number
of occupied children. Retain the most populous class and then the
same number of children at every surviving parent. Work from the
bottom level upwards; earlier equalities are preserved because
whole subtrees are retained. Multiplying the losses over the $m$
levels proves the estimate.
\end{proof}

\begin{remark}[The two parameters in uniformization]
\label{rem:uniform-parameters}
WZ2 Lemma 7.9 writes $\delta=2^{-\eta^{-1}T}$ but gives retention
$(2T)^{-\eta^{-1}}$, while Definition 7.8 selects scales separated
by $2^{-\eta^{-1}}$. The block length and the number of levels are
interchanged in these displays. The formulation above separates them:
there are $m$ levels, each of length $T$ in logarithmic scale, and
the loss is a power $m$ of the one-level factor. The formal
uniform-tree construction keeps these two integers separately.
\end{remark}

\begin{lemma}[A scale with good spacing; WZ2, Lemma 7.10]
\label{lem:spacing}
\lean{Kakeya.Assouad.ScaleProfileCrossingStatement,Kakeya.Assouad.scale_profile_crossing}
\uses{def:uniform-set,def:kt-frostman}
For every $\eps>0$, choose a sufficiently fine uniform scale grid.
If $A\subset[0,1)^n$ is uniform on this grid and
$\cE_\delta(A)\geq\delta^{-s}$, then one can choose
$\rho>\delta^{1-\eps}$ such that each occupied $\rho$-cube, after
rescaling to unit size, has a
$(\delta/\rho,s,(\delta/\rho)^{-4n\eps})$-Frostman bound.
\end{lemma}
\begin{proof}
\uses{lem:uniform-set,prop:os-multiscale}
Decompose the logarithmic scale interval into pieces on which the
branching behaves like a fixed dimension $t_j$. Merge adjacent
pieces whose dimensions are in decreasing order. The resulting
dimensions increase, and their weighted average is at least
$s-\eps^2$. Hence the first piece with $t_j\geq s-\eps$ starts
before the terminal part of the scale interval. All subsequent pieces
have dimension at least $s-\eps$. Combining their bounds gives the
Frostman estimate in every surviving cube. The linked arithmetic
lemma makes the early-crossing argument explicit; the geometric
pruning and scale-selection results supply its hypotheses.
\end{proof}

\begin{lemma}[Replication of a parameter block; WZ2, Lemma 7.11]
\label{lem:block-replication}
\lean{Kakeya.Assouad.ParameterBlockAmplificationStatement,Kakeya.Assouad.parameter_block_amplification}
\uses{def:kt-frostman}
Let $0<s\leq n$ and $0<\delta\leq\rho\leq1$. Let $A_0$ lie in a
cube of side $\rho$ and be a
$(\delta,s,C)$-Katz--Tao set, with
$\cE_\delta(A_0)\gtrsim(\rho/\delta)^s$. One can place translated
copies in a bounded region to obtain a
$(\delta,s,O(C))$-Katz--Tao set $A$ with
$\cE_\delta(A)\gtrsim\delta^{-s}$. Each selected replica cell
contains one translated copy of $A_0$.
More generally, if the local lower bound is only
$\cE_\delta(A_0)\geq\kappa(\rho/\delta)^s$ with $0<\kappa\leq1$,
the same construction gives
$\cE_\delta(A)\gtrsim\kappa\delta^{-s}$. The upper Katz--Tao
constant does not require $\kappa$ to be bounded away from zero.
\end{lemma}
\begin{proof}
Place about $\rho^{-s}$ copies on a grid of spacing comparable to
$\rho^{s/n}$. For a ball of radius at most $\rho$, use the original
Katz--Tao bound. Between $\rho$ and the grid spacing, only boundedly
many copies are met. Above the grid spacing, count the grid cells
and multiply by the size of one copy. Since $r^n\leq r^s$ for
$r\leq1$ and $s\leq n$, the desired bound follows.
\end{proof}

\begin{remark}[The replication cells]
\label{rem:replication}
The proof of WZ2 Lemma 7.11 prints $\rho^{s/n}\leq\rho$ (with the
dimension also denoted by $d$ there). The required inequality is
$\rho^{s/n}\geq\rho$ for $0<\rho\leq1$ and $s\leq n$, which
keeps the copies apart. Also, the assertion about translated copies
refers to the chosen cells of the construction, not to every translated
axis-parallel cube of side $\rho$: an arbitrarily shifted cube can
cut a copy in two. The formalization carries explicit replica
indices and disjointness data.
\end{remark}

\begin{remark}[The density deficit in the linked replication theorem]
\label{rem:replication-deficit}
The extraction in Lemma~\ref{lem:frostman-to-kt} loses a Frostman
constant and a logarithm, so its output need not satisfy a saturated
cardinality bound. This loss must survive replication.
The linked three-parameter Lean statement makes a more specific
choice than the general paper formulation: in a block of scale $q$,
it assumes an upper Katz--Tao exponent $k\leq1$ and a lower
cardinality exponent $0<s\leq k$. It produces $\asymp q^{-1}$
disjoint copies and a one-dimensional Katz--Tao set, with the exact
cardinality identity
\[
 \#A^*=\#\{\text{copies}\}\,\#A_0
       \gtrsim q^{-1}(q/\delta)^s
       =\delta^{-1}h^{1-s},\qquad h=\delta/q.
\]
The small factor $h^{1-s}$ is retained in the projected mass bound.
Replacing this by $\#A^*\gtrsim\delta^{-1}$ would discard a
hypothesis needed in the cinematic argument.
\end{remark}

\section{The one-step estimate and cinematic curves}

\begin{lemma}[Freezing a line parameter in a thin slab]
\label{lem:freeze-line-parameter}
\uses{def:twisted-projection}
Suppose $|c-c_0|\leq\delta$ and $|t|\leq1$. The area-preserving
target shear $(u,t)\mapsto(u-c_0t,t)$ sends the projected line
$\pi_f(\ell_{a,b,c,d})$ to within horizontal distance $\delta$ of
the graph of $g_{a,b,d}(t)=a+bf(t)+dtf(t)$.
\end{lemma}
\begin{proof}
After the shear the first coordinate is
$a+bf(t)+dtf(t)+(c-c_0)t$. The last term has absolute value at
most $\delta$, and the shear has determinant one.
\end{proof}

\begin{lemma}[From cinematic projected mass to area]
\label{lem:cinematic-area}
\uses{thm:cinematic-l32}
Fix the constants in Theorem~\ref{thm:cinematic-l32} and an output
loss $a>0$. At a sufficiently small scale $\delta$, let $G$ satisfy
that theorem's hypotheses and choose measurable
$E_g\subset N_{\lambda\delta}(\graph g)$ for each $g\in G$.
Writing $Z=\bigcup_{g\in G}E_g$ and
$M_{\rm proj}=\sum_{g\in G}|E_g|$, one has
\[
 M_{\rm proj}\leq |Z|^{1/3}\delta^{-a},\qquad
 |Z|\geq\delta^{3a}M_{\rm proj}^{\,3}.
\]
Here $M_{\rm proj}$ is planar projected mass, not three-dimensional
shaded tube volume.
\end{lemma}
\begin{proof}
Put $H=\sum_{g\in G}\one_{N_{\lambda\delta}(\graph g)}$.
H\"older's inequality and the cinematic estimate give
$M_{\rm proj}\leq\int_Z H\leq |Z|^{1/3}\|H\|_{3/2}
\leq |Z|^{1/3}\delta^{-a}$. Cube and rearrange.
\end{proof}

\begin{lemma}[Fine area versus coarse area; WZ2, Lemma 7.12]
\label{lem:projection-one-step}
\uses{def:twisted-projection,def:tube-frostman}
In the normalized setting of Theorem~\ref{thm:twisted-parameter},
and with the source loss chosen sufficiently small, there are
$\rho\in(\delta^{1-\eps^2},1)$ and a subshading $Y'$ of density
at least $\delta^{4\eta}$ such that
\[
 |\pi_f(E_{Y'})|\geq(\delta/\rho)^\eps
                         |N_\rho(\pi_f(E_{Y'}))|.
\]
\end{lemma}
\begin{proof}
\uses{lem:uniform-set,lem:spacing,lem:frostman-to-kt,
  lem:block-replication,lem:freeze-line-parameter,
  lem:wz1-cinematic-family,lem:cinematic-area}
First pigeonhole the incidence multiplicity along fibres of the
twisted projection. Uniformly refine that level set and pull it back
to a subshading. Keeping one multiplicity level relates the retained
target area to retained incidence mass; area retention alone would
not give shading density. Then select popular tubes.

In the bounded line chart, choose a $\delta$-slab in the parameter
$c$, retaining its mass and parameter-count losses. This is a finite
slab selection, rather than an assertion that an exact $c$-fibre of
a finite net is rich. Lemma~\ref{lem:freeze-line-parameter} removes
the common affine term with an $O(\delta)$ target error. Uniformize
the resulting triple parameters, select a block of side $q$, and
apply Frostman-to-Katz--Tao extraction at the rescaled mesh
$h=\delta/q$. Return the extracted points to the original block
before replicating them. Retain the cardinality deficit described in
Remark~\ref{rem:replication-deficit}.

The jet comparison in Lemma~\ref{lem:wz1-cinematic-family} converts
parameter separation and Katz--Tao counts into the corresponding
$C^2$ metric conditions. The replicated, transported shadings then
satisfy Lemma~\ref{lem:cinematic-area}. Translating parameters by
$(a_0,b_0,d_0)$ translates each horizontal image slice by
$a_0+b_0f(t)+d_0tf(t)$, so all copies have the same projected area.
Subadditivity bounds the replicated area by the number of copies
times the original block area; disjoint parameter copies need not
have disjoint projected images. Dividing by that number gives a
lower area bound for the block. Its curves lie in an $O(q)$
neighbourhood of one representative curve. The target uniformity
transfers this local comparison to the whole retained image.
All these steps use one compatible chain of subshadings.
\end{proof}

\begin{remark}[The three scales in the one-step argument]
\label{rem:projection-scales}
Write $q$ for the side of the selected parameter block and
$h=\delta/q$ for its mesh after unit rescaling. The coarse target
neighbourhood is at scale $r=q$, up to fixed geometric constants;
it is not at scale $h$. The conclusion above denotes this target
scale by $\rho$.
If $q>\delta^{1-\theta}$, then $\delta\leq h<\delta^\theta$.
In particular,
\[
                  \delta^{-a}\leq h^{-a/\theta}.
\]
A cinematic loss measured in $\delta$ therefore needs its own
budget: to spend at most $h^{-b}$, choose $a\leq\theta b$.
The absolute $C^2$ bound, neighbourhood dilation, and counting
constants are fixed before choosing the cinematic scale threshold.
Extraction, shading, and fixed-constant losses consume additional
parts of this budget.
\end{remark}

\begin{theorem}[Cinematic $L^{3/2}$ estimate; WZ2, Theorem 7.13]
\label{thm:cinematic-l32}
\lean{Kakeya.Cinematic.WZ2UniformC2Input,Kakeya.Cinematic.wz2_uniformC2_input}
\leanok
\uses{def:cinematic-family}
Fix the cinematic and doubling constants, a uniform absolute $C^2$
bound $M$, a Katz--Tao constant $C_0$, and a neighbourhood dilation
$\lambda\geq1$. For every $\eps>0$, and all sufficiently small
$\delta$, a $\delta$-separated finite subfamily $G$ satisfying
\[
       \#(G\cap B(g,r))\leq C_0(r/\delta)
       \quad(r\geq\delta)
\]
obeys
\[
       \left\|\sum_{g\in G}\one_{N_{\lambda\delta}(\graph g)}
       \right\|_{L^{3/2}(\R^2)}\leq\delta^{-\eps}.
\]
All the fixed constants, including $M$, precede the choice of the
small-scale threshold.
\end{theorem}
\begin{proof}
\uses{thm:pyz}
\leanok
Apply the cinematic maximal theorem on a bounded collection of
coordinate strips, with an auxiliary loss smaller than the requested
$\eps$. Sum the strip estimates, pass from the integral of the
$3/2$ power to the $L^{3/2}$ norm, and absorb the fixed constants
by decreasing the scale threshold.
\end{proof}

\begin{remark}[The normalized analytic formulation]
\label{rem:cinematic-normalization}
The paper states Theorem 7.13 for a fixed cinematic family and allows
a small-power loss in the counting constant. The Lean interface
above fixes a Katz--Tao constant, a graph-neighbourhood dilation, and
an absolute $C^2$ bound before choosing the fine scale. This is the
form used after selecting and rescaling the parameter blocks. The
need to keep the absolute bound in a family-uniform version of the
PYZ theorem is explained in Remark~\ref{rem:pyz-uniformity}.
\end{remark}

\section{Coarsening and the finite iteration}

\begin{lemma}[Uniform parameter fibres preserve the normalized count]
\label{lem:uniform-parameter-coarsening}
\uses{def:kt-frostman,def:uniform-set}
Let $A$ be a nonempty finite indexed family of line parameters in
$\R^4$ satisfying
\[
 \#\{p\in A:\|p-x\|_\infty\leq u\}\leq Cu^2\#A
       \quad(x\in\R^4,\ u\geq\delta).
\]
Let $\delta\leq q\leq1$. Partition $A$ into fibres $A_v$, indexed
by a finite parent family $P$, with
$\|p-v\|_\infty\leq q$ for $p\in A_v$. Suppose
$m\leq\#A_v\leq Km$ for every parent, where $m>0$ and $K\geq1$.
Then
\[
 \#\{v\in P:\|v-x\|_\infty\leq r\}
                       \leq4KC r^2\#P\quad(r\geq q).
\]
\end{lemma}
\begin{proof}
Each counted parent contributes at least $m$ distinct fine indices
within distance $2r$ of $x$. Their total is at most $4Cr^2\#A$,
while $\#A\leq Km\#P$. Cancel $m$.
\end{proof}

In the projection argument this lemma is applied to a selected
parameter family whose normalized count has already been restored.
Uniform child counts supply the required fibre bounds. The coarse
shading is induced by thickening the retained fine shadings; its
density, height support, and compatibility with the target
neighbourhood are checked separately. These geometric properties,
together with the count, permit the next application of the one-step
estimate. The carrier containment predicate alone does not provide
this parameter-space conclusion; see
Remark~\ref{rem:parameter-frostman}.

\begin{lemma}[Finite scale growth and telescoping]
\label{lem:projection-iteration}
Let $0<\delta<1$ and $0<\theta<1$. Starting at $r_0=\delta$,
suppose that whenever $r_j<\delta^\theta$ a next scale can be chosen
with $r_j^{1-\theta}<r_{j+1}\leq1$. The first index $N$ with
$r_N\geq\delta^\theta$ satisfies
\[
 N\leq1+\left\lceil\frac{\log\theta}{\log(1-\theta)}\right\rceil.
\]
If positive areas $A_j$ obey
$A_j\geq(r_j/r_{j+1})^a A_{j+1}$ and
$A_N\geq c r_N^2$, with $a,c>0$, then
\[
 A_0\geq c(\delta/r_N)^a r_N^2
                       \geq c\delta^{a+2\theta}.
\]
\end{lemma}
\begin{proof}
Induction gives $r_j\geq\delta^{(1-\theta)^j}$ as long as the
process continues. Once $(1-\theta)^j\leq\theta$, the stopping
condition holds, giving the stated finite bound. Multiplying the
area inequalities cancels every intermediate scale. Finally use
$r_N^{-a}\geq1$ and $r_N^2\geq\delta^{2\theta}$.
\end{proof}

\begin{remark}[The terminal scale and the loss schedule]
\label{rem:projection-iteration}
For the paper's iteration, the terminal area is the area of an
$r_N$-neighbourhood of a nonempty retained image. It contains a disk
of radius $r_N$, giving the quadratic terminal bound in the lemma.
The product of ratios ends at $r_N$.
Choose a small one-step exponent $a>0$ and use its available scale
gain: for example, $\theta=a^2$ in Lemma~\ref{lem:projection-one-step}.
Arrange $a+2\theta<\eps$, then fix the finite number of calls before
choosing the initial density and nonconcentration losses. If each
call spends an additional factor $\delta^b$, reserve $Nb$ in the
exponent and require $a+2\theta+Nb<\eps$. A final decrease of the
scale threshold absorbs fixed constants.
WZ2's displayed iteration count and last ratio are not needed
literally: scale growth gives the finite bound above, and telescoping
uses the actual terminal index. The Lean implementation keeps a
finite scale schedule, transition data, and a terminal estimate.
\end{remark}

\section{The original WZ1 slope and norm estimates}

For comparison, we record the statements in WZ1 Sections 6--7.
Extremality in this section has the WZ1 meaning of
Definition~\ref{def:wz1-extremal}. These statements use direction
separation and $L^{3/2}$ norms; WZ2 replaces them by the CWA and
projected-area statements above.

\begin{lemma}[A dense tube segment with small projection; WZ1, Lemma 6.2]
\label{lem:wz1-dense-segment}
\uses{def:ad-set,def:tube}
Let $0\leq\alpha<1$, $0<\delta\leq\rho\leq1$, and $\eps>0$.
Let $U$ be a segment of length $\sqrt\rho$ in a $\delta$-tube $T$,
with its full transverse thickness, and suppose
$F\subset U$ satisfies $|F|\geq\delta^\eps|U|$. If $v\in S^2$
and $v\cdot F$ is a $(\rho,\alpha,\delta^{-\eps})_1$-ADset, then
\[
                 |v\cdot\dir(T)|
                     \lesssim\delta^{-2\eps/(1-\alpha)}\sqrt\rho.
\]
\end{lemma}
\begin{proof}
A longitudinally dense subset of $U$ has a projection whose length
is controlled below by the angle of its axis to $v^\perp$.
The AD covering bound at scale $\rho$ bounds that projection above.
Comparing the two estimates and raising to the power $1/(1-\alpha)$
gives the conclusion.
\end{proof}

\begin{lemma}[Large derivative on a populated interval; WZ1, Lemma 6.3]
\label{lem:wz1-large-derivative}
\uses{prop:wz1-c2-grains}
For every $\eps>0$, choose $\eta$ and $\delta$ sufficiently small.
Suppose a WZ1 smooth-grain configuration with loss $\eta$ has
$\delta$-separated directions. If an interval $J\subset[-1,1]$
satisfies $\sqrt\delta\leq|J|\leq\delta^\eps$,
\[
 \sum_{T\in\T}|Y(T)\cap(\R^2\times J)|\geq\delta^\eta|J|,
\]
and $E_Y\cap(\R^2\times J)$ is covered by at most
$\delta^{-\eta}|J|^{-2+\sigma_3}$ cubes of side $|J|$, then
$|f'(z)|\geq|J|$ throughout $J$.
\end{lemma}
\begin{proof}
\uses{lem:wz1-dense-segment}
If the derivative is small at one point, the second derivative bound
keeps the global grain direction nearly constant across $J$. The
dense-segment lemma then confines many directions near one plane,
inconsistent with the direction count and the mass in the interval.
This is the directional predecessor of WZ2 Lemma 6.6.
\end{proof}

\begin{lemma}[A slab with comparable derivative; WZ1, Lemma 6.4]
\label{lem:wz1-comparable-derivative}
\uses{prop:wz1-c2-grains}
For every $\eps>0$ and scale ceiling, there is a WZ1 smooth-grain
configuration below that ceiling for which all shadings lie in
$[-1,1]^2\times[0,\delta^\eps]$, and a number $m\geq\delta^\eps$
such that
\[
             m\leq |f'(z)|\leq2m\quad(0\leq z\leq\delta^\eps).
\]
\end{lemma}
\begin{proof}
\uses{lem:wz1-large-derivative}
Choose a height interval retaining sufficient shaded mass and the
extremal spatial covering bound. Obtain the derivative lower bound
there, then localize and translate so that its magnitude is comparable
throughout one retained slab. Start with a smaller source loss to
retain the requested grain and extremality bounds.
\end{proof}

\begin{proposition}[Nonsingular WZ1 grains; WZ1, Proposition 6.1]
\label{prop:wz1-nonsingular}
\uses{def:wz1-extremal,def:plane-map,def:ad-set}
At arbitrarily small scales there are WZ1 $\eps$-extremal
configurations with global grains whose slope obeys
$1\leq|f'|\leq2$ and $|f''|\leq1/100$, together with local grains
and a $1$-Lipschitz plane map. The two AD covering bounds are those
of Proposition~\ref{prop:wz1-grains}, with the same retained shading.
The plane map is evaluated on that shaded union, as in
Remark~\ref{rem:plane-domain}.
\end{proposition}
\begin{proof}
\uses{lem:wz1-comparable-derivative,lem:mild-rescaling}
Rescale the selected height slab and the transverse coordinates.
The first derivative becomes comparable to one, and choosing the
slab sufficiently short makes the normalized second derivative small.
Transport the plane map and the local covering estimates with the
same affine change of variables.
\end{proof}

\begin{corollary}[A large projection norm; WZ1, Corollary 6.5]
\label{cor:wz1-large-norm}
\uses{def:twisted-projection,def:wz1-extremal}
For every $\eps>0$ and scale ceiling, there are a WZ1
$\eps$-extremal tube configuration below that ceiling and a slope
with $1\leq|f'|\leq2$, $|f''|\leq1/100$, such that
\[
 \left\|\sum_{T\in\T}\one_{\pi_f(T)}\right\|_{L^{3/2}(\R^2)}
                         \geq\delta^{-1-\sigma_3/3+\eps}.
\]
\end{corollary}
\begin{proof}
\uses{prop:wz1-nonsingular}
The global grains make the projected shaded union small, while the
sum of the projected shaded masses remains large. H\"older's
inequality bounds the $L^{3/2}$ norm below by this mass divided by
the one-third power of the support area. Enlarging each projected
shading to $\pi_f(T)$ can only increase the norm.
\end{proof}

\begin{proposition}[A small projection norm; WZ1, Proposition 7.1, normalized slope]
\label{prop:wz1-small-norm}
\uses{def:twisted-projection}
For every $\eps>0$, and all sufficiently small $\delta$, a family
of $\delta$-tubes with $\delta$-separated directions in WZ1's
bounded chart satisfies
\[
 \left\|\sum_{T\in\T}\one_{\pi_f(T)}\right\|_{L^{3/2}(\R^2)}
                           \leq\delta^{-1-\eps}
\]
whenever $f(0)=0$, $1\leq|f'|\leq2$, and $|f''|\leq1/100$.
\end{proposition}
\begin{proof}
\uses{thm:wz1-cinematic,lem:wz1-cinematic-family}
Split the four line parameters into approximately $\delta^{-1}$
windows of one slope parameter $c$. In each window, subtracting
$ct$ leaves a cinematic family. Direction separation supplies its
one-dimensional counting bound. Apply the cinematic estimate and
sum over the windows.
\end{proof}

\begin{remark}[An absolute normalization of the slope]
\label{rem:wz1-slope-offset}
The displayed paper statement of Proposition 7.1 specifies derivative
bounds but no bound on $f(0)$. Its uniform conclusion needs an
absolute normalization, which is present in the later cinematic
family construction and in the formal twisted-projection input.
Indeed, take a single vertical tube and $f(z)=A+z$. The derivative
conditions hold for every $A$, but the projected tube has area
comparable to $A\delta$ for large positive $A$, so its indicator has
$L^{3/2}$ norm comparable to $(A\delta)^{2/3}$. At a fixed scale this
is unbounded as $A\to\infty$. A fixed bound on $f(0)$, or the
normalization $f(0)=0$ used above, removes this obstruction.
\end{remark}

\begin{theorem}[The WZ1 cinematic estimate; WZ1, Theorem 7.2]
\label{thm:wz1-cinematic}
\uses{def:cinematic-family}
Let $\cF\subset C^2(I)$ be a fixed cinematic family on a compact
interval. For every $\eps>0$, at sufficiently small $\delta$, a
$\delta$-separated finite $F\subset\cF$ satisfying
\[
              \#(F\cap B)\leq\delta^{-\eps}(r/\delta)
\]
for every ball $B$ of diameter $r\geq\delta$ obeys
\[
               \left\|\sum_{g\in F}\one_{N_\delta(\graph g)}
                    \right\|_{3/2}\leq\delta^{-\eps}.
\]
For a threshold uniform over families, an absolute $C^2$ bound must
be fixed first, as explained in Remark~\ref{rem:pyz-uniformity}.
The completed Lean interface with a fixed counting constant is
Theorem~\ref{thm:cinematic-l32}.
\end{theorem}
\begin{proof}
\uses{thm:pyz}
Use the one-dimensional counting case of the cinematic maximal
estimate, localize the coordinate domain, and absorb fixed constants
with the fine-scale threshold.
\end{proof}

\begin{lemma}[The three-parameter cinematic family; WZ1, Lemma 7.3]
\label{lem:wz1-cinematic-family}
\uses{def:cinematic-family}
If $f(0)=0$, $1\leq|f'|\leq2$, and $|f''|\leq1/100$ on
$[0,1]$, the functions
\[
          g_{a,b,d}(t)=a+bf(t)+dtf(t),\qquad a,b,d\in[-1,1],
\]
form a cinematic family with constants independent of $f$.
They also have a uniform absolute $C^2$ bound.
\end{lemma}
\begin{proof}
The determinant of the two-jets of $1,f,tf$ is
$2(f')^2-ff''$, which is uniformly bounded away from zero because
$|f|\leq2$. Thus the jet at any point controls all three parameter
differences. This proves nondegeneracy and the equivalence of the
parameter and $C^2$ metrics, from which doubling follows.
\end{proof}

